%==============================================================================
% Chapitre 29 - Outils mécaniques pour la balistique
%==============================================================================

\section{Introduction}

Les chapitres précédents ont introduit plusieurs notions importantes :

\begin{itemize}
    \item masse ;
    \item centre de gravité ;
    \item centre de pression ;
    \item moments d'inertie.
\end{itemize}

Afin de comprendre les phénomènes de stabilité, il est nécessaire
d'introduire quelques outils mécaniques supplémentaires.

L'objectif de ce chapitre n'est pas de développer un cours complet de
mécanique, mais de fournir les notions indispensables à la compréhension des
chapitres suivants.

\begin{aretenir}

Les outils présentés dans ce chapitre seront utilisés pour expliquer la
stabilité, la précession et la nutation des projectiles.

\end{aretenir}

\section{Grandeurs scalaires et vectorielles}

Certaines grandeurs physiques sont entièrement décrites par une valeur
numérique.

On les appelle des scalaires.

Exemples :

\begin{itemize}
    \item masse ;
    \item température ;
    \item pression ;
    \item énergie.
\end{itemize}

D'autres nécessitent également une direction.

On les appelle des vecteurs.

Exemples :

\begin{itemize}
    \item vitesse ;
    \item accélération ;
    \item force ;
    \item déplacement.
\end{itemize}

\section{Notion de vecteur}

Un vecteur possède :

\begin{itemize}
    \item une norme ;
    \item une direction ;
    \item un sens.
\end{itemize}

Graphiquement, il est souvent représenté par une flèche.

\begin{exemple}

Une vitesse de :

\begin{equation}
800~m/s
\end{equation}

vers le nord est différente d'une vitesse de :

\begin{equation}
800~m/s
\end{equation}

vers l'est.

\end{exemple}

\section{Addition de vecteurs}

Les vecteurs peuvent être additionnés.

Cette propriété permet notamment d'expliquer :

\begin{itemize}
    \item les tirs depuis un véhicule ;
    \item l'effet du vent ;
    \item certains changements de référentiel.
\end{itemize}

\begin{exemple}

Un projectile quittant un véhicule en mouvement possède :

\begin{itemize}
    \item sa vitesse propre ;
    \item la vitesse du véhicule.
\end{itemize}

La vitesse observée résulte de l'addition de ces deux vecteurs.

\end{exemple}

\section{La force}

Une force est une action mécanique capable de modifier le mouvement d'un
corps.

Elle est représentée par un vecteur.

Ses caractéristiques sont :

\begin{itemize}
    \item intensité ;
    \item direction ;
    \item sens ;
    \item point d'application.
\end{itemize}

\section{Rappel de la deuxième loi de Newton}

Nous avons déjà rencontré :

\begin{equation}
F = ma
\end{equation}

Cette équation signifie que la force produit une accélération.

Plus la masse est importante, plus il faut appliquer une force élevée pour
obtenir la même accélération.

\section{Le bras de levier}

Considérons une porte.

Il est plus facile de l'ouvrir en poussant sur la poignée qu'au voisinage des
charnières.

Cette différence s'explique par le bras de levier.

Le bras de levier correspond à la distance entre :

\begin{itemize}
    \item l'axe de rotation ;
    \item le point d'application de la force.
\end{itemize}

\section{Moment d'une force}

Une force appliquée à distance d'un axe produit un moment.

Le moment mesure la capacité de cette force à provoquer une rotation.

Dans le cas simple :

\begin{equation}
M = F \times d
\end{equation}

où :

\begin{itemize}
    \item $M$ représente le moment ;
    \item $F$ représente la force ;
    \item $d$ représente le bras de levier.
\end{itemize}

\begin{aretenir}

Une même force produit un effet plus important lorsqu'elle agit loin de
l'axe de rotation.

\end{aretenir}

\section{Application aux projectiles}

Lorsqu'une force aérodynamique agit au centre de pression alors que le centre
de gravité est situé ailleurs, un moment apparaît.

Ce moment tend à faire pivoter le projectile.

C'est l'un des mécanismes fondamentaux de la stabilité.

\section{Rotation}

Un projectile rayé tourne autour de son axe longitudinal.

Cette rotation est caractérisée par une vitesse angulaire.

La vitesse angulaire indique la rapidité de la rotation.

\section{Vitesse angulaire}

La vitesse angulaire est généralement notée :

\begin{equation}
\omega
\end{equation}

Elle s'exprime en :

\begin{equation}
rad/s
\end{equation}

Plus la vitesse angulaire est élevée, plus le projectile tourne rapidement.

\section{Tour par minute}

Dans certains domaines, on utilise également :

\begin{equation}
rpm
\end{equation}

pour \emph{revolutions per minute}.

Les projectiles modernes peuvent atteindre plusieurs centaines de milliers de
tours par minute.

\begin{exemple}

Un projectile quittant un canon rayé peut dépasser :

\begin{equation}
200\,000~rpm
\end{equation}

voire davantage selon le calibre et le pas de rayure.

\end{exemple}

\section{Quantité de mouvement angulaire}

Un corps en rotation possède une quantité de mouvement angulaire.

Cette grandeur joue pour la rotation un rôle analogue à celui de la quantité
de mouvement en translation.

Plus elle est élevée, plus il est difficile de modifier l'orientation du
corps.

\section{Effet gyroscopique}

Une toupie en rotation reste droite beaucoup plus facilement qu'une toupie
immobile.

Ce phénomène est appelé effet gyroscopique.

Le même principe intervient pour les projectiles rayés.

La rotation contribue à leur stabilité.

\section{Perturbations}

Pendant son vol, un projectile subit de nombreuses perturbations :

\begin{itemize}
    \item vent ;
    \item turbulences ;
    \item défauts géométriques ;
    \item variations aérodynamiques.
\end{itemize}

La manière dont il réagit dépend notamment :

\begin{itemize}
    \item de sa rotation ;
    \item de ses moments d'inertie ;
    \item des moments aérodynamiques appliqués.
\end{itemize}

\section{Importance pour la balistique}

Les notions introduites dans ce chapitre permettent de comprendre :

\begin{itemize}
    \item la stabilité statique ;
    \item la stabilité gyroscopique ;
    \item la précession ;
    \item la nutation ;
    \item le spin drift.
\end{itemize}

Elles constituent donc des outils essentiels pour l'étude avancée des
projectiles.

\section{Vers la stabilité statique}

Nous disposons désormais de tous les éléments nécessaires pour répondre à une
question fondamentale :

\begin{quote}
Comment savoir si un projectile tend naturellement à revenir vers sa position
d'équilibre ou à s'en éloigner ?
\end{quote}

La réponse conduit à la notion de stabilité statique.

\section{À retenir}

\begin{aretenir}

\begin{itemize}
    \item Une force est une grandeur vectorielle.
    \item Une force appliquée à distance crée un moment.
    \item Le bras de levier influence ce moment.
    \item Un projectile rayé possède une vitesse angulaire élevée.
    \item La rotation produit un effet gyroscopique stabilisateur.
\end{itemize}

\end{aretenir}
